\section{TRIT: ternary orientation, carry and quadratic representation}
\label{trit:capitulo}

The construction of APP jointly preserves position, the
additive and multiplicative evaluations and their nonadic decompositions. The oriented ternary reduction
now acts on these data. Its task is twofold: to distinguish
the three operational classes of an evaluation and to preserve the integer coordinate
that allows it to be restored. This distinction is subsequently involved in the phase
selection of the topological phase kernel, called TPK after its initials in
English, \emph{Topological Phase Kernel}.

The term \emph{trit} designates a ternary unit. In the balanced convention
of HMT its values are \(-1,0,+1\). The same alphabet admits different
functions depending on the domain: reading an arithmetic evaluation, selecting
an operational mode or indexing a quadratic realization. Each function
is defined with its domain and the information it preserves.

The development begins with the integer evaluations of APP. Quadratic
realizations appear later as representations of the tritic type already
determined. The concrete construction of their generators and the Euler
extension are developed in chapter~\ref{eul:capitulo}, after introducing
the calendar and the TPK mode selection.

\subsection{Balanced ternary reduction of APP evaluations}
\label{trit:reduccion}

At a position \((i,j)\), with \(i,j\in\{1,\ldots,9\}\), the integer
evaluations are \(S(i,j)=i+j\) and \(P(i,j)=ij\). Each preserves its positive
remainder modulo nine and its quotient. Ternary reduction is applied to the
corresponding evaluation and produces another pair of coordinates.

\begin{definicion}[Oriented ternary decomposition]
\label{trit:def:balanceada}
Let
\[
 \mathcal T=\{-1,0,+1\},\qquad
 c_3(n)=\left\lfloor\frac{n+1}{3}\right\rfloor,\qquad
 r_3(n)=n-3c_3(n),\quad n\in\mathbb Z.
\]
The decomposition and restitution maps are
\[
 \mathcal B_3:\mathbb Z\longrightarrow\mathcal T\times\mathbb Z,
 \qquad \mathcal B_3(n)=(r_3(n),c_3(n)),
\]
\[
 \mathcal L_3:\mathcal T\times\mathbb Z\longrightarrow\mathbb Z,
 \qquad \mathcal L_3(r,c)=r+3c.
\]
The first component is the balanced ternary representative; the second
is the extension quotient.
\end{definicion}

\begin{proposicion}[Exact restitution and uniqueness]
\label{trit:prop:restitucion}
The maps \(\mathcal B_3\) and \(\mathcal L_3\) are inverses.
In particular, for every integer \(n\),
\[
 n=r_3(n)+3c_3(n),\qquad r_3(n)\in\mathcal T,
\]
and this expression is unique.
\end{proposicion}
\begin{proof}
If \(c=\lfloor(n+1)/3\rfloor\), then
\(3c\le n+1<3c+3\). Therefore \(-1\le n-3c<2\), and the integer
\(r=n-3c\) belongs to \(\mathcal T\). If
\(r+3c=s+3d\), with \(r,s\in\mathcal T\), the integer \(r-s\) is a multiple
of three and lies between \(-2\) and \(2\). Thus, \(r=s\) and \(c=d\).
Finally, for any pair \((r,c)\),
\(\lfloor(r+3c+1)/3\rfloor=c\), since \(0\le r+1\le2\).
This proves both inverse compositions.
\end{proof}

The oriented section of the quotient \(\mathbb Z/3\mathbb Z\) is fixed by
\[
 [0]_3\longmapsto0,\qquad [1]_3\longmapsto+1,
 \qquad [2]_3\longmapsto-1.
\]
Replacing \(2\) with \(-1\) is accompanied by the necessary quotient:
\(2=-1+3\). The integer evaluation remains recoverable in the pair.

\begin{tablacompleta}
\caption{Balanced decomposition of the nine positions on an APP line.}
\label{trit:tabla:nueve}
\begin{tabular}{c*{9}{r}}
\toprule
\(n\)&1&2&3&4&5&6&7&8&9\\
\midrule
\(r_3(n)\)&1&$-1$&0&1&$-1$&0&1&$-1$&0\\
\(c_3(n)\)&0&1&1&1&2&2&2&3&3\\
\bottomrule
\end{tabular}
\notatabla{Coincidence of ternary representatives preserves three
distinct positions in each class; their quotients restore the integers.}
\end{tablacompleta}

For example, at \((i,j)=(5,5)\), the sum equals \(10\) and the product equals
\(25\). Both ternary readings give \(+1\), while their
complete decompositions are
\[
 \mathcal B_3(10)=(1,3),\qquad
 \mathcal B_3(25)=(1,8).
\]
The common signature identifies a residue type; the sheet and quotient distinguish
the two evaluations. The APP state also preserves the position that
produced them.

\subsubsection{Compatibility between moduli nine and three}

Write the nonadic decomposition of an evaluation as
\[
 n=R_9(n)+9q_9(n),\qquad R_9(n)\in\{1,\ldots,9\}.
\]
Reducing its positive representative ternarily yields
\[
 R_9(n)=r_3(R_9(n))+3c_3(R_9(n)).
\]
Substituting,
\begin{equation}
 r_3(n)=r_3(R_9(n)),\qquad
 c_3(n)=c_3(R_9(n))+3q_9(n).
\label{trit:eq:compatibilidad39}
\end{equation}
Uniqueness of the balanced decomposition proves both identities.
The visible reduction modulo three can be performed on the nonadic
remainder, while the complete reconstruction also transports the
previous quotient.

In the product \(5\cdot5=25\), the nonadic decomposition is
\(25=7+9\cdot2\). Since \(7=1+3\cdot2\), the preceding formula gives
\(25=1+3(2+3\cdot2)=1+3\cdot8\). Both reading orders restore
the same integer because the quotients compose explicitly.

\subsection{Arithmetic of ternary states with carry}
\label{trit:aritmetica}

Preserving the pairs makes it possible to transfer the APP operations to
balanced coordinates. The transfer determines the additive carry and
the cross terms of the product.

\begin{definicion}[Balanced additive carry]
For \(r,s\in\mathcal T\), let
\[
 r\oplus s=r_3(r+s),\qquad
 \chi(r,s)=\frac{r+s-r_3(r+s)}3.
\]
The integer \(\chi(r,s)\) is the carry produced when normalizing the sum of
the two representatives.
\end{definicion}

\begin{tablacompleta}
\caption{Additive carry of the balanced ternary section.}
\begin{tabular}{c|rrr}
\toprule
\(\chi(r,s)\)&\(s=-1\)&\(s=0\)&\(s=+1\)\\
\midrule
\(r=-1\)&$-1$&0&0\\
\(r=0\)&0&0&0\\
\(r=+1\)&0&0&1\\
\bottomrule
\end{tabular}
\notatabla{The two extremes encode
\(-1-1=1-3\) and \(1+1=-1+3\).}
\end{tablacompleta}

\begin{proposicion}[Recoverable operations]
\label{trit:prop:operaciones}
If \(n=r+3c\) and \(m=s+3d\), with \(r,s\in\mathcal T\), then
\begin{align}
 (r,c)\boxplus(s,d)
 &=\bigl(r\oplus s,c+d+\chi(r,s)\bigr),
 \label{trit:eq:suma}\\
 (r,c)\boxtimes(s,d)
 &=\bigl(rs,rd+sc+3cd\bigr)
 \label{trit:eq:producto}
\end{align}
are, respectively, the decompositions of \(n+m\) and \(nm\).
\end{proposicion}
\begin{proof}
The sum decomposes as
\[
 n+m=r+s+3(c+d)
     =r_3(r+s)+3\bigl(c+d+\chi(r,s)\bigr).
\]
In the product,
\[
 nm=(r+3c)(s+3d)=rs+3(rd+sc+3cd).
\]
The integer \(rs\) belongs to \(\mathcal T\). Both expressions are
normalized; uniqueness in proposition~\ref{trit:prop:restitucion}
identifies their components.
\end{proof}

These laws distinguish accumulation and multiplication in the record.
The sum aggregates the quotients and adds the residue carry. The product
also combines each representative with the quotient of the other evaluation.
This arithmetic content is maintained when the operations are incorporated
into a TPK transition.

\begin{ejemplo}[Sum and product of two evaluations]
From \(2=-1+3\) and
\(4=1+3\) follows
\[
 (-1,1)\boxplus(1,1)=(0,2),\qquad
 (-1,1)\boxtimes(1,1)=(-1,3).
\]
Restitution produces \(6\) and \(8\), respectively. In the second
operation, the quotient is \((-1)\cdot1+1\cdot1+3\cdot1\cdot1=3\).
\end{ejemplo}

\begin{proposicion}[Cocycle identity and associativity]
\label{trit:prop:cociclo}
For \(r,s,t\in\mathcal T\),
\begin{equation}
 \chi(r,s)+\chi(r\oplus s,t)
 =\chi(s,t)+\chi(r,s\oplus t).
\label{trit:eq:cociclo}
\end{equation}
The operations \(\boxplus\) and \(\boxtimes\), together with their identity
elements and additive inverses transported by \(\mathcal B_3\), make
\(\mathcal T\times\mathbb Z\) a ring isomorphic to \(\mathbb Z\).
\end{proposicion}
\begin{proof}
When normalizing \((r+s)+t\), the total quotient is the first side of
\eqref{trit:eq:cociclo}; when normalizing \(r+(s+t)\), it is the second.
The final representatives coincide because they represent the same class
modulo three. Uniqueness of the complete decomposition requires the
quotients to coincide as well. By proposition~\ref{trit:prop:operaciones},
\(\mathcal L_3\) is a bijection preserving sum and product. It therefore transports
associativity, distributivity and commutativity of the
integer operations. Its identity elements are \((0,0)\) and \((1,0)\).
\end{proof}

The cocycle proof expresses a concrete compatibility of the record:
normalization can be performed in stages without changing the final
evaluation. If the sequence of operations is of interest, the record also preserves
the order of the normalizations performed.

\subsection{Ternary refinement and recovery at arbitrary depth}
\label{trit:refinamiento}

The decomposition is iterated on the quotient. Given \(n_0\in\mathbb Z\),
define recursively
\[
 r_k=r_3(n_k),\qquad n_{k+1}=c_3(n_k),\qquad k\ge0.
\]
Each transition satisfies \(n_k=r_k+3n_{k+1}\). Composing these
equalities produces a restitution at any prescribed depth.

\begin{proposicion}[Restitution of a truncation]
\label{trit:prop:profundidad}
For every \(N\ge1\),
\begin{equation}
 n_0=\sum_{k=0}^{N-1}3^kr_k+3^Nn_N.
\label{trit:eq:restitucionN}
\end{equation}
The sequence \((r_0,\ldots,r_{N-1})\), together with the terminal quotient
\(n_N\), determines \(n_0\) exactly.
\end{proposicion}
\begin{proof}
For \(N=1\), the identity is the initial decomposition. If it holds for
\(N\), substitute \(n_N=r_N+3n_{N+1}\) into the last term to
obtain the formula for \(N+1\). Restitution also reconstructs the
intermediate quotients, starting at \(n_N\) and applying
\(n_k=r_k+3n_{k+1}\) in reverse order.
\end{proof}

\begin{ejemplo}[A carry crossing three levels]
The evaluation \(26\) successively gives
\[
 26=-1+3\cdot9,\qquad 9=0+3\cdot3,
 \qquad 3=0+3\cdot1,\qquad1=1+3\cdot0.
\]
Its representatives, ordered from lower to higher power, are
\((-1,0,0,1)\). Restitution is \(26=-1+27\). The zero symbol at the
intermediate levels preserves a position and a determined quotient.
\end{ejemplo}

\begin{proposicion}[Termination for a fixed integer]
Iteration of the balanced quotient reaches zero for every initial
integer. Once trailing zeros are removed, the finite balanced ternary
expansion of an integer is unique.
\end{proposicion}
\begin{proof}
For \(|n|\ge2\),
\( |c_3(n)|\le(|n|+1)/3<|n|\).
For \(n\in\{-1,0,1\}\), the quotient is zero. The absolute value
strictly decreases while it is at least two, so the
iteration terminates. If two expansions represent the same integer,
their reduction modulo three fixes the same first representative; subtracting it
and dividing by three repeats the argument for all positions.
\end{proof}

This result concerns the arithmetic of a fixed integer. In
TPK extensions, the record length increases and new
evaluations can be generated. There, depth corresponds to a compatible
sequence of states, whose construction requires the survival
rules. Formula~\eqref{trit:eq:restitucionN} provides the local
restitution mechanism that those extensions must preserve.

\subsection{Two ternary coordinates of a nonadic class}
\label{trit:nonadico}

The reduction of \(\mathbb Z/9\mathbb Z\) modulo three has three elements
in each fiber. The additional coordinate can be written as a residue
ternary quotient.

\begin{proposicion}[Nonadic coordinates with carried addition]
\label{trit:prop:beta}
The map
\[
 \beta:\mathcal T\times\mathbb F_3\longrightarrow\mathbb Z/9\mathbb Z,
 \qquad \beta(r,\bar c)=[r+3c]_9
\]
is a bijection of sets. Nonadic addition and multiplication take the
form
\begin{align}
 \beta(r,\bar c)+\beta(s,\bar d)
 &=\beta\bigl(r\oplus s,\overline{c+d+\chi(r,s)}\bigr),
 \label{trit:eq:beta-suma}\\
 \beta(r,\bar c)\beta(s,\bar d)
 &=\beta\bigl(rs,\overline{rd+sc}\bigr).
 \label{trit:eq:beta-producto}
\end{align}
\end{proposicion}
\begin{proof}
Changing \(c\) to \(c+3k\) changes \(r+3c\) by \(9k\), so
the map is well defined. If two images coincide, reduction
modulo three fixes the same \(r\). The remaining difference, divided by
three, fixes the same class \(\bar c\). Injectivity between two sets
of nine elements proves bijectivity. The operations come from
\eqref{trit:eq:suma} and \eqref{trit:eq:producto}; the term \(3cd\) in the
multiplicative quotient vanishes upon reduction modulo three.
\end{proof}

The carry establishes the composition law between the two levels. For
example,
\[
 \beta(1,0)+\beta(1,0)=\beta(-1,1).
\]
The second coordinate records the step performed by the integer sum before
its residue reading. In particular, \(\mathbb Z/9\mathbb Z\) has
elements of order nine, whereas the additive group of
\(\mathbb F_3^2\) has exponent three. The preceding bijection preserves
the sum through the explicit carried law.

The three APP positions \(3,6,9\) are written
\[
 [3]_9=\beta(0,1),\qquad
 [6]_9=\beta(0,2),\qquad
 [9]_9=\beta(0,0).
\]
On the ideal \(3\mathbb Z/9\mathbb Z\), extraction of the second
coordinate is
\[
 [3c]_9\longmapsto[c]_3.
\]
This is an additive bijection with \(\mathbb F_3\). It differs from
direct reduction modulo three, which assigns zero to all three elements.
Moreover, the product of two elements of this ideal is zero modulo nine.
Its multiplicative structure thus retains a function different from
multiplication in the field \(\mathbb F_3\).

\subsubsection{The fiber of six ternary positions}

An oriented word of length six belongs to \(\mathcal T^6\),
with \(3^6=729\) possibilities. The choice of representatives induces a
bijection with \(\mathbb F_3^6\). Grouping positions in pairs and
applying \(\beta\) also yields a set-theoretic bijection with
\((\mathbb Z/9\mathbb Z)^3\). The transported law uses the carry
of each pair.

The census of words actually emitted by the TPK is performed after
constructing its emitter and chronology. The ambient fiber of \(729\)
elements fixes coordinate capacity; the set of emissions and its ternary
image are determined by the operational rules. This separation makes it possible
to preserve, in the subsequent census, both multiplicities and the
distinct records sharing a visible word.

\subsection{Oriented involution and transport memory}
\label{trit:involucion}

\begin{proposicion}[Compatibility with sign reversal]
For all \(n\in\mathbb Z\),
\[
 \mathcal B_3(-n)=(-r_3(n),-c_3(n)).
\]
\end{proposicion}
\begin{proof}
From \(n=r+3c\) we deduce \(-n=(-r)+3(-c)\). Since \(-r\in\mathcal T\),
uniqueness of the decomposition proves the statement.
\end{proof}

For a sequence of signatures of length \(m\), define
\[
 \mathcal C_{\rm rev}(\tau_1,\ldots,\tau_m)
 =(-\tau_m,\ldots,-\tau_1).
\]
The second application restores the initial order and signs, so
that \(\mathcal C_{\rm rev}^{2}=\operatorname{id}\). In a sequence
of pairs, the same rule simultaneously transports
\((r_k,c_k)\mapsto(-r_{m+1-k},-c_{m+1-k})\). The length and each
restorable evaluation are determined.

Division memory and transport memory serve different
purposes. The pair \((r_3(n),c_3(n))\) reconstructs an integer evaluation.
Reconstructing a trajectory also requires the initial position,
orientations, sequence of modes and boundary crossings. Two paths
can reach the same evaluation and preserve different transport
records. The enriched state integrates both orders of information.

In the temporal interpretation adopted by HMT, \(+1\) corresponds to
return with memory, \(0\) to the present threshold and \(-1\) to the opening of
conjugate extensions. This assignment fixes a regime reader.
Effective evolution requires the TPK transition on the complete
records; its iterations determine which phase returns and which memory is
incremented.

\subsection{Quadratic algebras of the three regimes}
\label{trit:cuadraticas}

The tritic index has an algebraic realization through a quadratic
relation. Its integral formulation precedes the choice of the scalar
field and allows discrete information to be separated from subsequent
analytic realizations.

\begin{definicion}[Quadratic algebra of tritic type]
For \(\tau\in\mathcal T\), define
\[
 \mathcal Q_{\tau,\mathbb Z}=\mathbb Z[u]/(u^2+\tau).
\]
For a commutative ring \(A\), its realization by change of scalars is
\[
 \mathcal Q_{\tau,A}=A[u]/(u^2+\tau).
\]
The class of \(u\), written \(J_\tau\), satisfies
\(J_\tau^2=-\tau I\).
\end{definicion}

Division by the monic polynomial \(u^2+\tau\) provides a unique
expression \(aI+bJ_\tau\). The product is calculated by reducing only the
quadratic term:
\begin{equation}
 (aI+bJ_\tau)(cI+dJ_\tau)
 =(ac-\tau bd)I+(ad+bc)J_\tau.
\label{trit:eq:producto-cuadratico}
\end{equation}
This multiplication represents the quadratic regime. The arithmetic of
balanced pairs, for its part, preserves integers and their quotients
through \eqref{trit:eq:suma} and \eqref{trit:eq:producto}.

\subsubsection{Nine-element realizations}

Over \(\mathbb F_3\), the three algebras have nine elements and admit
the following classification:
\[
 \mathcal Q_{+1,\mathbb F_3}\cong\mathbb F_9,\qquad
 \mathcal Q_{0,\mathbb F_3}=\mathbb F_3[\varepsilon]/(\varepsilon^2),
 \qquad
 \mathcal Q_{-1,\mathbb F_3}\cong\mathbb F_3\times\mathbb F_3.
\]

\begin{proposicion}[Distinction of the three finite types]
\label{trit:prop:tres-tipos}
The first algebra is a field; the second contains a nonzero
nilpotent; the third has two nontrivial orthogonal idempotents whose
sum is the identity. The three are pairwise nonisomorphic.
\end{proposicion}
\begin{proof}
The polynomial \(u^2+1\) takes the values \(1,2,2\) on \(0,1,2\);
it is therefore irreducible over \(\mathbb F_3\). Its quotient is a
field of degree two. In the second case, the class \(\varepsilon\) is
nonzero and satisfies \(\varepsilon^2=0\). In the third,
\(u^2-1=(u-1)(u+1)\), with coprime factors. The evaluation
\(a+bu\mapsto(a+b,a-b)\) is an isomorphism to
\(\mathbb F_3\times\mathbb F_3\). The elements
\(2(1+u)\) and \(2(1-u)\) are its two coordinate idempotents.
These algebraic properties distinguish the three types.
\end{proof}

In the development of exceptional incidence, a concretely
constructed ternary operator satisfying \(A_W^2=-I\) induces the action of
\(\mathbb F_9\) on its space. If that space is \(\mathbb F_3^6\),
the dimension over \(\mathbb F_9\) is three: the action is defined by
\((a+b\iota)v=av+bA_Wv\). The operator fixes multiplication by the
quadratic unit. The corresponding Paley construction supplies that
operator in its causal place; the present classification establishes the
content of the relation it must satisfy.

\subsubsection{Conjugation, algebraic norm and inverse}

In the real realization, conjugation and the algebraic norm are
\[
 \overline{aI+bJ_\tau}=aI-bJ_\tau,\qquad
 N_\tau(aI+bJ_\tau)=a^2+\tau b^2.
\]

\begin{proposicion}[Multiplicative norm and inverse restitution]
\label{trit:prop:norma}
For \(z,w\in\mathcal Q_{\tau,\mathbb R}\),
\[
 z\bar z=N_\tau(z)I,
 \qquad N_\tau(zw)=N_\tau(z)N_\tau(w).
\]
An element \(z\) is invertible exactly when \(N_\tau(z)\ne0\);
in that case \(z^{-1}=\bar z/N_\tau(z)\).
\end{proposicion}
\begin{proof}
The product \eqref{trit:eq:producto-cuadratico} gives
\((aI+bJ_\tau)(aI-bJ_\tau)=(a^2+\tau b^2)I\).
Conjugation is multiplicative and the algebra is commutative, so
\((zw)\overline{zw}=(z\bar z)(w\bar w)\). If the norm is nonzero,
the formula exhibits an inverse. If \(zw=I\), multiplicativity gives
\(N_\tau(z)N_\tau(w)=1\), forcing \(N_\tau(z)\ne0\).
\end{proof}

The matrix representation
\[
 aI+bJ_\tau\longmapsto
 \begin{pmatrix}a&-\tau b\\b&a\end{pmatrix}
\]
is faithful and has determinant \(a^2+\tau b^2\). For \(\tau=+1\), the
form is positive definite. For \(\tau=0\), it depends only on
\(a\), while \(J_0\) retains a nonzero linear component. For
\(\tau=-1\), the form has signature \((1,1)\); the directions
\(I\pm J_{-1}\) are zero divisors.

In particular, the word \emph{norm} here designates a multiplicative quadratic
form. The positive norms of Hilbert spaces, used
in other realizations of the state, retain their own definition.

\subsubsection{Composition of ratios between coordinates}

A direction with nonzero scalar component is represented by
\(I+uJ_\tau\). Multiplication produces
\[
 (I+uJ_\tau)(I+vJ_\tau)
 =(1-\tau uv)I+(u+v)J_\tau.
\]
When \(1-\tau uv\ne0\), the ratio between the components of the product is
\begin{equation}
 u\star_\tau v=\frac{u+v}{1-\tau uv}.
\label{trit:eq:razones}
\end{equation}
The scalar factor \(1-\tau uv\) forms part of the complete product.
If only the ratio is retained, that factor must be recorded when
it is involved in a subsequent reconstruction. If the denominator vanishes,
the product is examined in its two components: it can represent a
direction with zero scalar component or the zero element.

The quadratic classification thus prepares a precise reading of the three
regimes. Selection of a regime preserves its index, arithmetic
preserves its carries and transport preserves the oriented sequence.
The TPK composes these data through explicit operators. On that
composition will be built the concrete generators and the preservation
of forms that articulates the Euler extension.
